\chapter{The Interpretive Problem}
\label{chap:interpretation-problem}

\section{Semantic Nakedness}

When humans interpret the output of a computer monitor, they rely on an astronomical stack of shared conventions (encoding standards,
pixel geometry, scan order, handedness, row length, etc.). Changing these conventions can render the identical underlying data unrecognizable or meaningless.

A raw informational structure does not announce what it represents. 

The exact same finite set of bits can be interpreted as a number, a computation, an entire universe, random noise, or nothing at all.


\subsection*{Base}

Consider Alice using a decimal system, motivated by biological convenience (ten fingers):
\[
2 + 3 = 5.
\]

Bob, working in a computational environment, adopts binary representation:
\begin{itemize}
\item $2 \rightarrow 10_2$
\item $3 \rightarrow 11_2$
\item $5 \rightarrow 101_2$
\end{itemize}

In binary arithmetic the same relation becomes:
\[
10_2 + 11_2 = 101_2.
\]

If Alice encountered Bob's notation without context, she would misinterpret the symbols.
In particular, $10_2$ does not refer to ten.
Even the full expression requires an assumed encoding scheme.
Thus, meaning is not intrinsic to the symbol structure alone, but depends on an interpretive convention.

\subsection*{Symbolic Relativity}

The arbitrariness of representation becomes clearer when we replace the digits entirely.
The binary alphabet $\{0,1\}$ may be replaced with coin faces $\{H,T\}$:
\[
HT + HH = HTH.
\]

Operational symbols such as $+$ and $-$ are likewise part of a chosen syntactic convention, not intrinsic to the underlying relation.

\subsection*{Two's Complement}

In finite digital systems, identical bit patterns admit multiple semantic interpretations.
For an 8-bit word:
\begin{itemize}
\item $11111111$ may represent $255$ in unsigned interpretation,
\item or $-1$ in two's complement representation.
\end{itemize}

The physical state is identical; only the interpretation differs.

\subsection*{Reading Order}

Even ordering conventions are not intrinsic.
A bitstring such as $001$ may correspond to different values depending on bit significance:
\begin{itemize}
\item left-to-right (MSB first): $1$,
\item right-to-left (LSB first): $4$.
\end{itemize}

These are merely two of the $n!$ possible permutations of bit significance.

\subsection*{Code as Data}

Consider a computer in full physical detail.
A memory dump yields only a finite sequence of bits.
The CPU, instruction set, and storage devices are themselves physically instantiated information structures.
At the level of raw physical state, there is no intrinsic separation between ``code'' and ``data''.
This distinction arises solely through an interpretive layer.

Computation, or any other semantics, arises only from relations between states under a chosen interpretive rule.


\section{Interpretive Explosion}

A raw bitstring does not intrinsically contain logic, syntax, semantics, truth conditions, or consistency relations.
These arise only after selecting an interpretive scheme - with the presense of an interpreter (observer).

As the number of bits increases, the space of possible interpretations grows combinatorially.
A single finite configuration admits a vast number of syntactically valid decompositions into code/data partitions,
memory layouts, instruction interpretations, and execution histories.

While the state space itself contains only $2^n$ configurations, the number of possible transition functions is $(2^n)^{2^n}$.
For $n=3$ bits this is already $8^8 = 16\,777\,216$.
For $n=10$ the number has thousands of digits.



\section{Equivalence Classes - Interpretive Gauge Invariance}

Even after an interpretation is assigned, it is not unique.
Different interpretive schemes may preserve the same underlying relational structure, as demonstrated
in Chapter~\nameref{chap:the-problem-with-time}.

Physics has long grappled with this exact representational ambiguity, demonstrating that distinct mathematical formulations can map to identical physical realities.
In mathematical physics, when a physical state remains unchanged under a change of local descriptive frameworks, we call that symmetry a gauge invariance.

Consider a sufficiently large bitstring.
Under one interpretation it may appear as random noise.
Under another it may be interpreted as implementing the relation $2 + 3 = 5$, as encoding a mathematical structure, or even as the informational substrate of an observer.
Instruction layouts, memory partitions, addressing schemes, and representational conventions may all vary, yet certain structural relations can remain invariant across these transformations.

We therefore define:

\begin{tcolorbox}[title=Interpretive Equivalence Class]
Two interpretations belong to the same interpretive equivalence class if they preserve the same observable relational structure.
\end{tcolorbox}

What ultimately matters is not the precise arrangement of symbols, but the invariant relational structure that survives reinterpretation:
\[
\text{information} \ \longrightarrow \ \text{interpretation} \ \longrightarrow \ \text{relational structure (invariant)}
\]



\section{The Observer as Information}

In static, timeless universe interpretation is not an external reading imposed on a bitstring of the universe, but
an internal relation between informational substructures.

The interpretive rule space grows enormously faster than state space. 

The DNA simulation in Chapter~\nameref{chap:running-dna-simulation} traces out just one single axiomatic description of Alice, a subset, of what is called the Alice equivalence class. 
The number of ways the information in the simulation can be interpreted as Alice is astronomically larger. Vast majority of them do not correspond to any recognizable computation, 
mathematics.


\begin{tikzpicture}[scale=1.0, font=\small]
  % Unknown: the outside (assumption: drawn as the surrounding region)
  \draw[thick, dashed, gray, fill=gray!5, rounded corners=8pt] (-4.7,-4.7) rectangle (4.7,5.5);
  \node[font=\small\bfseries, gray!70!black] at (0,5.0) {Unknown};

  % Circles tangent at the bottom (y = -4.2) so each layer has a top band for its label
  \draw[thick, fill=blue!5, draw=blue!40]     (0,0)     circle (4.2cm);
  \draw[thick, fill=orange!10, draw=orange!50, dashed] (0,-0.9)  circle (3.3cm);
  \draw[thick, fill=teal!15, draw=teal!60]    (0,-1.8)  circle (2.4cm);
  \draw[thick, fill=green!15, draw=teal!70]   (0,-2.7)  circle (1.5cm);

  \node[align=center, font=\small\bfseries, blue!80!black]   at (0,3.3)  {Static Information ($\mathbf{U}$)\\\footnotesize $2^n$ states (semantically silent)};
  \node[align=center, font=\small\bfseries, orange!80!black] at (0,1.5)  {Interpretive Landscape\\\footnotesize combinatorial explosion of mappings};
  \node[align=center, font=\small\bfseries, teal!80!black]   at (0,-0.3) {Mathematical Slice\\\footnotesize gauge orbits \& invariant relations};
  \node[align=center, font=\small\bfseries, teal!80!black]   at (0,-2.7) {Alice\\\footnotesize equivalence class};

  % Filters: each crosses exactly one boundary, at the circle's left edge
  \draw[-{Stealth}, thick, gray] (-3.9,-0.9) -- (-3.0,-0.9) node[midway, above, font=\tiny, align=center] {Information\\filter};
  \draw[-{Stealth}, thick, gray] (-2.9,-1.8) -- (-2.0,-1.8) node[midway, above, font=\tiny, align=center] {Mathematics\\filter};
  \draw[-{Stealth}, thick, gray] (-2.0,-2.7) -- (-1.1,-2.7) node[midway, above, font=\tiny, align=center] {Observer\\filter};
\end{tikzpicture}



\section{The Core Thesis}

By synthesizing the semantic nakedness of raw data with the mathematical necessity of gauge invariance, we arrive at the thesis:

\begin{principle}{Self-Interpretive Information Principle (SIIP)}
{siip}
The raw informational substrate is semantically silent.

Within this silent substrate exists a vast landscape of possible interpretive frameworks, equivalence classes, and emergent self-running observers.
\end{principle}


Or perhaps:

\begin{principle}{Internal-Emergence Principle (IEP)}
{iep}

In static universe, observers emerge as built-in interpreters, running all possible interpretations in parallel. Those interpretations that are self-aware, 
find themselve within the largest equivalence class of them. 

\end{principle}




\section{The Deep Mystery of Mathematics}

The preceding analysis allows us to revisit and resolve Eugene Wigner's famous puzzle concerning the "unreasonable effectiveness of mathematics in the natural sciences."

The effectiveness is an emergent consequence of data compression.

We find ourselves within an equivalence class of minimal description, because that dominates the measure  maximizing our probability of existence.

The shortest possible description of observers is one in which observers consist of the same informational substructures, a concept well known in software engineering.

We are all made of the same fundamental building blocks.

Consequently, mathematics is unreasonably effective because the interpreters themselves (observers) are highly unified.

This perspective dissolves Wigner's paradox. The universe does not need to be inherently mathematical; it only needs to be rich enough in raw information to allow compressed structures to form.

Mathematics is the ultimate emergent compession language, not a fundamental substrate of the universe.

Mathematics works so well because we are highly unified filters.

Invariant relational structure is fundamental; mathematics is the observer-discovered description of that structure.



\subsection*{Reality and Physics}

Suppose we have two observers, Alice and Bob.

\begin{itemize}
 \item Alice has an internally realized informational structure \(O_A\).
 \item Bob has an internally realized informational structure \(O_B\).
 \item Each has their own observer walk and subjective experience (dreams, thoughts, feelings etc.)
 \item Some of their observations, interpretations, and informational structures overlap.
\end{itemize}

We can express the idea schematically:

\[
R\_{AB}=O_A\cap O_B
\]

Here, the intersection represents their shared informational structure.

With more observers:

\[
R=\bigcap\_{i=1}^{N}O_i
\]


\begin{itemize}
\item Subjective reality: What an individual observer internally experiences.
\item Objective reality: The stable, reproducible intersection of many independent observer perspectives.
\end{itemize}



In a self-contained, static universe, everything an observer sees must already exist inside their own information structure. 
Alice contains Bob's structure, and Bob contains Alice's. Reality is the invariant equivalence class born from this mutual symmetry.

Mathematics is an emergent compression engine operating inside the observer's mind. It allows Alice to re-map her internal data in astronomically many ways:

Imagination/Dreaming: Transformations that describe configurations outside her local trajectory.

Reality (Planning/Prediction): Transformations whose ordinal mappings match the physical configurations she subsequently traverses.

Math feels objective and universal because observers share identical low-level sub-structures. 
Just as software engineering uses modular code to save resources, the underlying dynamics of the system favor minimal description length. 
Alice and Bob reason with the same mathematics because they are built from the same shared structural primitives.

In static timeless universe, whatever Alice observes as Bob, must already be encoded into the structure of Alice herself. Vast majority of 
the interpretations of the information in Alice describe Alice-like structures, namely Bob, and six billion others that we observe.

As a compression language, mathematics allow also descrptions that Alice cannot observe in her structure. This is what Alice calls her 
dreams, thoughts, and imagination.  

Some of those reinterpretations (compressed descriptions) may actually be part of her structure, she will only discover it later (when her 
walk through configurations coincide). This is what Alice calls "planning the future". 

For example, two arrangements giving the same Alice, like $2+2+2$ and $2\cdot3$, is the equivalence class. 
A mapping Alice computes about her “future” counts as reality if the ordinal configurations she later passes through match it, 
and as imagination if they don't. 

